% Transcription — fonds Alexandre Grothendieck, shelfmark 10, batch 3 (pages 41-60).
% Established from the facsimile archives/batches/10/batch-03.pdf (a mirror of
% https://grothendieck.umontpellier.fr/10.pdf), per the skill
% transcribe-grothendieck.
%
% Pass: Opus 5 (claude-opus-5), 2026-09-11 — first pass, unchecked against the
% pages by a human.
%
% The batch holds two typescripts, neither of them his, both corrected in his
% hand.
%
% Pages 41-57 continue the typescript begun page 32 and broken off mid-word at
% the foot of page 40. Page 41 closes n° 10.16; page 42 opens « 11. —
% Préschémas en groupes affines », which runs to the last written line of the
% batch:
%
%   42-44 the affine algebra A(X) = f_*(O_X); Lemme 11.1 on when
%         A(X) ⊗ A(Y) → A(X ×_S Y) is an isomorphism, with a margin note of his
%         calling the formula false; Corollaire 11.2, the functor
%         X ↝ Spec A(X) commuting with finite products; Définition 11.3 of the
%         affine envelope G_af; 11.4 on Hom_{O_S}(V(F), V(E)); Proposition 11.5
%         and its proof, where his hand restores Lazard's theorem in its true
%         form — every flat module is a filtered colimit of free modules of
%         finite type — against a typed statement that had it backwards.
%   45-46 11.6, the two axioms of an operation of G on E written both
%         geometrically (two diagrams in θ) and coalgebraically ((CM1), (CM2)
%         in ρ), and the bijection between operations of G and of G_af; the
%         induced operation on a stable submodule; Lemme 11.8 for A free with
%         basis (φ_i); Proposition 11.9.
%   47-51 Corollaire 11.10 (E is the filtered union of its finite-type stable
%         submodules) with Remarque 11.10.1 crediting J.P. Serre and replacing
%         a reference to a seminar by the Publ. Math. paper; then Proposition
%         11.11 — affine ⟺ quasi-affine ⟺ faithful action on a quasi-affine
%         scheme ⟺ on a vector space ⟺ closed subgroup of some Gl(n) — first
%         drafted page 47 and cancelled there, restated page 48 with five
%         conditions, proved through Aut(V_i), the projective limit G' = lim G_i
%         and Lemme 11.11.1; Remarque 11.11.1 records what Raynaud showed over
%         a regular base of dimension ≤ 1; Lemme 11.12 and Proposition 11.13,
%         G as a projective limit of affine groups of finite type with
%         faithfully flat transition morphisms.
%   52-57 Lemme 11.14 (u faithfully flat ⟺ u° injective); Définitions 11.15,
%         semi-invariant vectors, their weight and associated character — the
%         word « poids » moved by hand onto « invariant » and its old
%         definition struck; Théorème 11.16 (Chevalley), H as the stabiliser of
%         a line of semi-invariants in Λ^d V; Théorème 11.17 (Chevalley), the
%         fpqc quotient G/H is affine, proved first for G algebraic and then in
%         general by limits; and Corollaire 11.18 on subgroups of a
%         diagonalisable group, cancelled by loops.
%
% Page 55 is the one leaf of the batch entirely in his hand: it supplies the
% algebraically closed case of 11.17 that page 56 asks for by a large circled
% « VOIR CI CONTRE » set against the struck reference to Chevalley. It runs
% through (Ker ρ)_red = H, then G/H affine by G/G° and G°/H ∩ G°; three N.B.
% follow, cancelled by diagonals, one of them noting that the reference to
% « BIBLE » is a corollary of 11.16.
%
% Pages 59-60 are a second typescript, on other paper and from another
% machine, and are the only leaves of the batch on the subject of the folder:
% a tannakian category over k of diagonalisable band D(Γ) is classified by
% ξ ∈ H²(k,G) ≃ Hom(Γ, Br(k)), i.e. by u : Γ → Br(k); the category is
% Γ-graded, the homogeneous objects are isotypic, d(i) is the order of u(i),
% K(M) ⊂ Z[Γ] is the set of Σ n_i e^i with n_i ≡ 0 mod d(i), and End(M_i) is
% the division algebra of invariant u(i). Page 60 reconstructs M from the
% Azumaya algebras A_i by a sub-additive category on N^(I), its Karoubian
% envelope, and the inversion of the L_i = det E_i, with two margin notes of
% his — « attention à contrainte de commutativité » and « A expliciter ! ».
%
% The letter M of those two pages is inked by hand throughout, the machine
% not having it. Insertions are \add{}, deletions \struck{}, his margin notes
% \marginal{}; typing instructions to the typist (« pas à la ligne ») are
% omitted. The cancelled runs of pages 49-51 and the crosses of page 50 are
% recorded where they fall.
%
% Page 58, the blank verso of page 59 where only its frappe shows through, is
% absent.
%
% The dating is the inventory's own, for the whole shelfmark.
\documentclass[11pt,a4paper]{article}
% Le préambule commun à toutes les transcriptions du fonds Grothendieck.
%
% Il tient en une idée : **l'incertitude se compose, elle ne se résout pas.**
% Une transcription qui lisse un mot douteux a détruit la seule chose pour
% laquelle on la faisait — un lecteur doit pouvoir savoir, à chaque endroit,
% ce qui était sur le feuillet et ce qui a été ajouté. Les six macros qui
% suivent sont donc l'appareil critique, et non de la décoration.
%
% Les mêmes commandes sont reconnues par `scripts/render.mjs`, qui produit la
% vue de lecture du site. Une macro ajoutée ici doit l'être aussi là-bas,
% faute de quoi le rendu échouera bruyamment — ce qui est voulu : un rendu
% silencieusement faux est pire qu'un rendu absent.

\ProvidesPackage{grothendieck}[2026/08/08 Transcriptions du fonds Grothendieck]

\usepackage{amsmath,amssymb,amsthm}
\usepackage{mathrsfs}
\usepackage{stmaryrd}
% Les diagrammes commutatifs sont partout dans ce fonds : c'est la langue
% même de la géométrie algébrique de Grothendieck, et une transcription qui
% les rendrait en prose ne transcrirait rien.
\usepackage{tikz-cd}
% Français par défaut — c'est la langue des feuillets — mais l'anglais est
% chargé aussi : la traduction et le résumé sont des documents anglais, et la
% typographie française (espaces avant les deux-points) y serait une faute.
% Ces documents basculent par \selectlanguage{english} en tête de corps.
\usepackage[english,french]{babel}
\usepackage{fontspec}
\usepackage{geometry}
\geometry{a4paper,margin=25mm,marginparwidth=28mm}
\usepackage{marginnote}
\usepackage{xcolor}
% ulem, not soul. `soul`'s \st and \ul are built on a character-by-character
% scan that XeLaTeX's font machinery breaks: the document fails with an
% undefined control sequence rather than degrading. ulem does the same two jobs
% and is Unicode-engine safe. `normalem` stops it hijacking \emph, which the
% transcriptions use for Grothendieck's own emphasis.
\usepackage[normalem]{ulem}
\usepackage{hyperref}
\hypersetup{colorlinks=true,linkcolor=black,urlcolor=black,pdfborder={0 0 0}}

% --- Métadonnées du lot -----------------------------------------------------
% Reprises telles quelles par le script de rendu, qui les affiche en tête de
% la vue de lecture. Elles fixent ce que le fichier prétend couvrir : sans
% elles, une transcription flotte sans point d'attache dans un fonds de seize
% mille feuillets.

\newcommand{\folder}[1]{\gdef\grothFolder{#1}}
\newcommand{\batch}[1]{\gdef\grothBatch{#1}}
\newcommand{\leaves}[2]{\gdef\grothFirst{#1}\gdef\grothLast{#2}}
\newcommand{\foldertitle}[1]{\gdef\grothFoldertitle{#1}}
\gdef\grothFolder{?}\gdef\grothBatch{?}\gdef\grothFirst{?}\gdef\grothLast{?}\gdef\grothFoldertitle{}

% --- Le résumé --------------------------------------------------------------
%
% Il ouvre la lecture modernisée, sous le titre « Résumé », et il est composé
% en petit. Ce n'est pas un ornement : c'est le seul passage du document qui
% ne soit pas des mathématiques, et un lecteur doit pouvoir le reconnaître
% avant de l'avoir lu — puis le sauter s'il connaît déjà le sujet, ou s'y
% arrêter s'il ne le connaît pas. Le filet qui le suit marque la reprise du
% corps.

\newenvironment{resume}
  {\small\setlength{\parskip}{0.45em}}
  {\par\medskip\hrule\medskip}

% --- Mots-clés ---------------------------------------------------------------
%
% En anglais, en clôture du résumé de la lecture modernisée : le vocabulaire
% moderne sous lequel chercher ce que les feuillets construisent. Le manifeste
% du site (`npm run manifest`) les extrait pour étiqueter la cote entière —
% cette ligne est la source unique des tags, il n'y a pas de fichier à côté.
\newcommand{\keywords}[1]{\par\smallskip\noindent
  {\footnotesize\textsc{Keywords} --- \textit{#1}}\par}

% --- L'appareil critique ----------------------------------------------------

% Le numéro de feuillet, en marge. C'est l'ancre qui permet de régler un
% désaccord sur une formule en regardant *un* feuillet plutôt qu'une chemise
% de six cents. Le site s'en sert aussi pour tourner les pages du fac-similé
% quand on fait défiler la transcription.
\newcommand{\leaf}[1]{%
  \marginnote{\footnotesize\color{gray!70}\textsf{#1}}%
  \label{leaf:#1}\ignorespaces}

% Illisible : on ne devine pas. Un mot inventé qui se lit comme les autres est
% le pire résultat possible.
\newcommand{\ill}{\textcolor{red!60!black}{[\,\ldots\,]}}

% Lecture douteuse : le mot est donné, et signalé comme incertain.
\newcommand{\uncertain}[1]{\uline{#1}}

% Ajout de l'éditeur — une lettre manquante, un mot sous-entendu.
\newcommand{\add}[1]{\textcolor{blue!50!black}{[#1]}}

% Biffé par Grothendieck lui-même. Ce qu'il a rayé fait partie du document :
% une rature dit souvent où la pensée a changé de direction.
\newcommand{\struck}[1]{\sout{#1}}

% Note du transcripteur, sur l'état matériel du feuillet.
\newcommand{\note}[1]{\par\smallskip{\small\color{gray!60!black}\textit{[#1]}}\par\smallskip}

% Note marginale de Grothendieck, distincte du corps du texte.
\newcommand{\marginal}[1]{\par\smallskip\noindent\hspace{1em}%
  {\small\color{gray!60!black}#1}\par\smallskip}

% --- Titre ------------------------------------------------------------------

\AtBeginDocument{%
  \begin{center}
    {\large\bfseries Fonds Alexandre Grothendieck --- cote n\textsuperscript{o}~\grothFolder}\\[2pt]
    {\itshape\grothFoldertitle}\\[4pt]
    {\small feuillets \grothFirst--\grothLast\ (lot \grothBatch)}\\[2pt]
    {\footnotesize\color{gray!60!black}Transcription établie d'après le fac-similé numérisé
     par l'Université de Montpellier.\\
     Les lectures incertaines sont soulignées, les illisibles notées [\,\ldots\,],
     les ajouts entre crochets.}
  \end{center}
  \vspace{1em}}

\endinput


\folder{10}
\batch{3}
\pages{41}{60}
\dating{[à partir de 1958]}
\watermark{Édition de démonstration}
\foldertitle{Catégories tensorielles : notes manuscrites (s.d.), tapuscrits (s.d.)}

\begin{document}

\section*{Préschémas en groupes affines (pages 41 à 57)}

\note{Suite du tapuscrit ouvert page 32 et interrompu au bas de la page 40 :
même machine, même papier, mêmes corrections à la main (crayon et encre
bleue). La page 41 achève le n\textsuperscript{o} 10.16 commencé au lot
précédent ; la page 42 ouvre un paragraphe nouveau, « 11. — Préschémas en
groupes affines », qui court jusqu'à la page 57. Dans tout ce paragraphe
$\mathcal{A}(X) = f_{*}(\mathcal{O}_{X})$, $\mathbb{V}$ et $\mathbb{W}$ sont
les foncteurs en modules attachés à un Module quasi-cohérent,
$\mathcal{H}om$ le foncteur des morphismes de $S$-foncteurs, et
$\mathrm{S}(\mathcal{F})$ l'Algèbre symétrique ; les lettres rondes de la
machine sont rendues par les rondes correspondantes. Ses insertions sont
données entre crochets, ses suppressions barrées, ses notes marginales en
retrait. La machine écorche « biunivoque » à plusieurs reprises et la main
le rétablit chaque fois ; ces corrections ne sont pas relevées une à une.}

\page{41}
\note{Fin du n\textsuperscript{o} 10.16, dont l'énoncé commence au bas de la
page 40 : « … soient $j \in I$, $G_{j}$ un $S_{j}$-grou- ».}

pe de présentation finie, $H_{j}$ un $S_{j}$-préschéma de présentation
finie à groupe d'opérateurs $G_{j}$. Pour que $H = H_{j} \times_{S_{j}} S$
soit un $S$-préschéma en espaces homogènes sous \struck{\ill{}} $G$ (resp.
un fibré principal homogène sous \struck{\ill{}} $G$) pour la
\emph{topologie fidèlement plate localement de présentation finie}, il faut
et il suffit qu'il existe un indice $i \in I$, $i \geqslant j$, tel que
$H_{i} = H_{j} \times_{S_{j}} S_{i}$ soit un $S_{i}$-préschéma en espaces
homogènes sous \struck{\ill{}} $G_{i}$ (resp. un fibré principal homogène
sous $G_{i}$).

\marginal{$G = G_{j} \times_{S_{j}} S$ \quad ; \quad $G_{i} = G_{j} \times_{S_{j}} S_{i}$}
\note{Les deux définitions marginales sont de sa main, accolées chacune à la
ligne où le symbole tapé a été biffé.}

Compte tenu de (10.14) et de (EGA IV 8.8.2, 8.8.3 et 8.10.5),
\struck{\ill{}} l'énoncé résulte de la propriété concernant les morphismes
\struck{\ill{}}\add{couvrants} pour la topologie fppf rappelée en (10.0).

\page{42}
\textbf{11. — Préschémas en groupes affines.}

\textbf{11.0.} — Notations : soient $S$ un préschéma, $X$ un $S$-préschéma,
$f : X \to S$ le morphisme structural ; on pose
$\mathcal{A}(X) = f_{*}(\mathcal{O}_{X})$.

\textbf{Lemme 11.1.} — Soient $X$ et $Y$ deux $S$-préschémas quasi-compacts
et quasi-séparés sur $S$, $f : X \to S$ et $g : Y \to S$ les morphismes
structuraux. Alors l'homomorphisme canonique
\[
\varphi : \mathcal{A}(X) \otimes_{\mathcal{O}_{S}} \mathcal{A}(Y)
\longrightarrow \mathcal{A}(X \times_{S} Y)
\]
est un isomorphisme dans chacun des cas suivants :
\begin{itemize}
\item[a)] $f$ et $g$ sont affines
\item[b)] $f$ ou $g$ est plat et affine
\item[c)] $g$ est plat et $f_{*}(\mathcal{O}_{X})$ est un
$\mathcal{O}_{S}$-Module plat.
\end{itemize}

\marginal{\uncertain{fausse} formule}
\marginal{$f_{*}(\mathcal{O}_{X})$}
\note{Deux notes de sa main, chacune entourée, en regard de la formule et de
la condition c). La première met en cause la formule affichée ; le premier
mot n'est pas assuré.}

Dans le cas b), on peut supposer que c'est $g$ qui est plat et affine.
Posons $S' = \operatorname{Spec} \mathcal{A}(X)$, $Y' = Y \times_{S} S'$,
$g' = g \times_{S} S'$ et notons $v$ le morphisme $S' \to S$. Alors $Y$ est
quasi-compact et quasi-séparé sur $S$, et, ou bien $S'$ est plat sur $S$
(cas c) d'après (EGA III (1.4.15.5)), ou bien $g$ est affine (cas a) et
b)). Alors (EGA III 1.4.15 et IV 1.2.7 ou II 1.5.2), on a :
\[
g'_{*}(\mathcal{O}_{Y'}) = v^{*} g_{*}(\mathcal{O}_{Y})
= \mathcal{A}(Y) \otimes_{\mathcal{O}_{S}} \mathcal{O}_{S'} .
\]

D'après (EGA II 1.2.7), $f$ se factorise à travers $v$ au moyen d'un
morphisme $p : G \to S'$, et on a : $X \times_{S} Y = X \times_{S'} Y'$ et
$p_{*}(\mathcal{O}_{X}) = \mathcal{O}_{S'}$. Puisque $v$ et $f$ sont
quasi-compacts et quasi-séparés, \struck{\ill{}} il en est de même de $p$
(EGA IV 1.2.2 et 1.2.4). \add{Ou} bien $Y$ est plat sur $S$ (cas b) et c))
et alors $Y'$ est plat sur $S'$, ou bien $X$ est affine sur $S$ (cas a)),
et alors $p$ est un isomorphisme. Dans le premier cas, on peut de nouveau
appliquer (EGA III 1.4.15 et IV 1.2.7), et dans les deux cas, on trouve :
\struck{\ill{}}
\[
p'_{*}(\mathcal{O}_{X \times_{S} Y}) = g'^{*} p_{*}(\mathcal{O}_{X})
= g'^{*}(\mathcal{O}_{S'}) = \mathcal{O}_{Y'} ,
\]
\add{où $p' = p \times_{S'} Y'$.}
\note{Le morphisme est tapé $p : G \to S'$ ; c'est $X$ que l'argument
demande. La définition de $p'$ est ajoutée de sa main en bout de ligne.}

D'où :
\[
\mathcal{A}(X \times_{S} Y) = v_{*} g'_{*} p'_{*}(\mathcal{O}_{X \times_{S} Y})
= v_{*} g'_{*}(\mathcal{O}_{Y'})
= v_{*}(\mathcal{A}(Y) \otimes_{\mathcal{O}_{S}} \mathcal{O}_{S'})
= \mathcal{A}(Y) \otimes_{\mathcal{O}_{S}} \mathcal{A}(X)
\]
d'après (EGA II 1.4.7).

\textbf{Corollaire 11.2.} — On définit un foncteur
$X \rightsquigarrow \operatorname{Spec} \mathcal{A}(X)$ de la
sous-catégorie pleine de $\mathrm{Sch}_{/S}$ des $X$-préschémas plats
quasi-compacts et quasi-séparés sur $S$ tels que $\mathcal{A}(X)$ soit un
$\mathcal{O}_{S}$-Module plat, dans celle des $S$-préschémas plats et
affines sur $S$. Ce foncteur commute aux produits finis, donc transforme
$S$-groupes en $S$-groupes.
\note{« des $X$-préschémas » est tapé ainsi, pour « des $S$-préschémas ».}

\page{43}
\textbf{Définition 11.3.} — Etant donné un $S$-groupe $G$ plat,
quasi-compact et quasi-séparé sur $S$, tel que $\mathcal{A}(G)$ soit plat
sur $\mathcal{O}_{S}$, nous noterons $G_{\mathrm{af}}$, et nous appellerons
\emph{enveloppe affine} de $G$ le $S$-groupe
$G_{\mathrm{af}} = \operatorname{Spec} \mathcal{A}(G)$.

\textbf{11.4.} — Etant donnés deux $\mathcal{O}_{S}$-Modules $\mathcal{E}$
et $\mathcal{F}$ quasi-cohérents, notons
$\mathcal{H}om_{S}(\mathbb{V}(\mathcal{F}), \mathbb{V}(\mathcal{E}))$ le
$S$-foncteur des morphismes de $S$-foncteurs de $\mathbb{V}(\mathcal{F})$
dans $\mathbb{V}(\mathcal{E})$, et
$\mathcal{H}om_{\mathcal{O}_{S}}(\mathbb{V}(\mathcal{F}), \mathbb{V}(\mathcal{E}))$
le sous-$S$-foncteur de
$\mathcal{H}om_{S}(\mathbb{V}(\mathcal{F}), \mathbb{V}(\mathcal{E}))$ formé
des $\mathcal{O}_{S}$-homomorphismes de $\mathcal{O}_{S}$-foncteurs en
modules (cf. Exp. I). Alors
$\mathcal{H}om_{\mathcal{O}_{S}}(\mathbb{V}(\mathcal{F}), \mathbb{V}(\mathcal{E}))(T)$
est en correspondance biunivoque canonique avec \struck{\ill{}}
$\operatorname{Hom}_{\mathcal{O}_{S}\text{-Mod.}}
(\mathcal{E} \otimes_{\mathcal{O}_{S}} \mathcal{O}_{T},
\mathcal{F} \otimes_{\mathcal{O}_{S}} \mathcal{O}_{T})$, \add{i.e.
$\mathcal{H}om_{\mathcal{O}_{S}}(\mathbb{V}(\mathcal{F}), \mathbb{V}(\mathcal{E}))
\simeq
\mathcal{H}om_{\mathcal{O}_{S}}(\mathbb{W}(\mathcal{E}), \mathbb{W}(\mathcal{F}))$.}

Soit $X$ un $S$-préschéma. Alors l'ensemble des morphismes de $S$-foncteurs
de $X$ dans
$\mathcal{H}om_{S}($\struck{\ill{}}$\mathbb{V}(\mathcal{F}), \mathbb{V}(\mathcal{E}))$
est en correspondance biunivoque avec
$\operatorname{Hom}_{S}(X \times_{S} \mathbb{V}(\mathcal{F}), \mathbb{V}(\mathcal{E}))$,
donc l'ensemble des morphismes de $S$-foncteurs de $X$ dans
$\mathcal{H}om_{\mathcal{O}_{S}}(\mathbb{V}(\mathcal{F}), \mathbb{V}(\mathcal{E}))$
est en correspondance biunivoque avec un sous-ensemble de
$\operatorname{Hom}_{S}(X \times_{S} \mathbb{V}(\mathcal{F}), \mathbb{V}(\mathcal{E}))
= \operatorname{Hom}_{\mathcal{O}_{S}\text{-Mod.}}
(\mathcal{E}, \mathcal{A}(X \times_{S} \mathbb{V}(\mathcal{F})))$
(EGA II 1.7.13).

\textbf{Proposition 11.5.} — Soient $X$ un $S$-préschéma quasi-compact
quasi-séparé sur $S$, $f : X \to S$ le morphisme structural, $\mathcal{E}$
et $\mathcal{F}$ deux $\mathcal{O}_{S}$-Modules quasi-cohérents. On suppose
vérifiée l'une des deux conditions suivantes :
\begin{itemize}
\item[a)] $f$ est affine
\item[b)] $\mathcal{F}$ est plat sur $\mathcal{O}_{S}$.
\end{itemize}
Alors :
\begin{itemize}
\item[i)] l'homomorphisme canonique
$\varphi : \mathcal{A}(X) \otimes_{\mathcal{O}_{S}} \mathrm{S}(\mathcal{F})
\longrightarrow \mathcal{A}(X \times_{S} \mathbb{V}(\mathcal{F}))$ est un
isomorphisme (où $\mathrm{S}(\mathcal{F})$ désigne l'Algèbre symétrique de
$\mathcal{F}$).
\item[ii)] l'ensemble des morphismes de $S$-foncteurs de $X$ dans
$\mathcal{H}om_{\mathcal{O}_{S}}(\mathbb{V}(\mathcal{F}), \mathbb{V}(\mathcal{E}))$
est en correspondance biunivoque avec
$\operatorname{Hom}_{\mathcal{O}_{S}\text{-Mod.}}
(\mathcal{E}, \mathcal{A}(X) \otimes \mathcal{F})$.
\end{itemize}

\marginal{$S$}
\note{Une lettre $S$ entourée, de sa main, en regard de i) — l'Algèbre
symétrique, que la machine note d'une $S$ ronde.}

Pour montrer i), d'après (11.1, a) et b)), il suffit de montrer que si
$\mathcal{F}$ est plat sur $\mathcal{O}_{S}$, \struck{\ill{}} alors
$\mathbb{V}(\mathcal{F})$ est plat sur $S$. D'après (EGA III 1.4.15.5),

\page{44}
il suffit de voir que si $\mathcal{F}$ est plat sur $\mathcal{O}_{S}$,
\add{il en est de même de $\mathrm{S}(\mathcal{F})$,} ce qui
\struck{résulte des définitions (EGA II 1.7.1), compte tenu d'}\add{est un
corollaire d'}un resultat dû à Daniel \textbf{Lazard}, suivant lequel
tout\struck{e limite inductive de modules plats est un} module plat
\add{sur un anneau est une limite inductive filtrante de modules libres de
type fini.}
\note{La correction rétablit le théorème de Lazard dans sa forme exacte :
la machine avait tapé « toute limite inductive de modules plats est un
module plat ».}

\struck{Remarquons de plus que
$\mathcal{A}(X) \otimes_{\mathcal{O}_{S}} \mathrm{S}\mathcal{F}
= \mathrm{S}(\mathcal{A}(X) \otimes_{\mathcal{O}_{S}} \mathcal{F})$
d'après (EGA II 1.1.2), ce qui montre que l'homomorphisme canonique
$\mathcal{A}(X) \otimes_{\mathcal{O}_{S}} \mathcal{F} \to
\mathcal{A}(X) \otimes_{\mathcal{O}_{S}} \mathrm{S}\mathcal{F}$ est
injectif.}
\note{Ces deux lignes sont biffées par une suite de diagonales et closes
d'une croix.}

Montrons maintenant ii). Soit $\alpha$ un morphisme de $S$-foncteurs de $X$
dans
$\mathcal{H}om_{\mathcal{O}_{S}}(\mathbb{V}(\mathcal{F}),$
$\mathbb{V}(\mathcal{E}))$ ;
il lui correspond un homomorphisme $\rho$ de $\mathcal{O}_{S}$-Modules de
$\mathcal{E}$ dans
$\mathcal{A}(X \times_{S} \mathbb{V}(\mathcal{F}))
= \mathcal{A}(X) \otimes_{\mathcal{O}_{S}} \mathrm{S}(\mathcal{F})$ tel que
si $\rho$ est une $S$-section de $X$, l'homomorphisme
\[
(\mathcal{A}(\eta) \otimes \mathrm{id}_{\mathrm{S}\mathcal{F}}) \circ \rho :
\mathcal{E} \longrightarrow
\mathcal{O}_{S} \otimes_{\mathcal{O}_{S}} \mathrm{S}\mathcal{F}
= \mathrm{S}\mathcal{F}
\]
corresponde à $\alpha(\eta)$.

\marginal{$\eta = ?$}
\note{Note de sa main, à l'encre bleue et entourée, en regard de cette
ligne : la $S$-section est tapée $\rho$, alors que c'est $\eta$ qu'elle
doit être — $\rho$ désigne déjà l'homomorphisme.}

Mais puisque $\alpha(\eta) \in{}$ \struck{\ill{}}
$\mathcal{H}om_{\mathcal{O}_{S}}(\mathbb{V}(\mathcal{F}), \mathbb{V}(\mathcal{E}))(S)$,
$\mathcal{A}(\eta) \otimes \mathrm{id}_{\mathrm{S}\mathcal{F}}$ se
factorise à travers l'homomorphisme canonique injectif
$\mathcal{F} \to \mathrm{S}\mathcal{F}$. On a donc, pour chaque
$S$-section $\eta$ de $X$, un diagramme commutatif :

\begin{tikzcd}[column sep=large]
\mathcal{E} \arrow[r, "\rho"] \arrow[d] &
  \mathcal{A}(X) \otimes \mathrm{S}\mathcal{F} \arrow[d]
  \arrow[dl, "\mathcal{A}(\eta) \otimes \mathrm{id}"] \\
\mathrm{S}\mathcal{F} & \mathcal{F} \arrow[l]
\end{tikzcd}

\note{Le diagramme est tracé à la main dans l'intervalle laissé par la
machine.}

ce qui montre que $\rho$ se factorise à travers l'homomorphisme canonique
injectif
$\mathcal{A}(X) \otimes_{\mathcal{O}_{S}} \mathcal{F} \longrightarrow
\mathcal{A}(X) \otimes_{\mathcal{O}_{S}} \mathrm{S}\mathcal{F}$, donc
$\rho$ peut être considéré comme un homomorphisme de
$\mathcal{O}_{S}$-Modules de $\mathcal{E}$ dans
$\mathcal{A}(X) \otimes_{\mathcal{O}_{S}} \mathcal{F}$.

Réciproquement, soit
$\rho \in \operatorname{Hom}_{\mathcal{O}_{S}\text{-Mod.}}
(\mathcal{E}, \mathcal{A}(X) \otimes_{\mathcal{O}_{S}} \mathcal{F})
\subset \operatorname{Hom}_{\mathcal{O}_{S}\text{-Mod}}
(\mathcal{E}, \mathcal{A}(X \times_{S} \mathbb{V}(\mathcal{F})))$ ; il lui
correspond un morphisme $\alpha$ de $X$ dans
$\mathcal{H}om_{S}(\mathbb{V}(\mathcal{F}), \mathbb{V}(\mathcal{E}))$.
Pour tout $S$-préschéma $T$ et toute $T$-section $\eta$ de
$X \times_{S} T$, $\alpha(\eta)$ correspond à
$\mathcal{A}(\eta) \otimes
\mathrm{id}_{\mathrm{S}\mathcal{F} \otimes \mathcal{O}_{T}}$,
homomorphisme de
$\mathcal{E} \otimes_{\mathcal{O}_{S}} \mathcal{O}_{T}$ dans
$\mathcal{F} \otimes_{\mathcal{O}_{S}} \mathcal{O}_{T}$, donc
$\alpha(\eta) \in
\mathcal{H}om_{\mathcal{O}_{S}}(\mathbb{V}(\mathcal{F}), \mathbb{V}(\mathcal{E}))(T)$,
ce qui montre que $\alpha$ est un morphisme de $X$ dans
$\mathcal{H}om_{\mathcal{O}_{S}}(\mathbb{V}(\mathcal{E}), \mathbb{V}(\mathcal{F}))$,
ce qui prouve ii).
\note{Le dernier $\mathcal{H}om$ est tapé avec ses deux arguments
intervertis.}

\struck{\ill{}}
\note{Les trois dernières lignes de la page — l'énoncé d'un résultat sur un
$S$-groupe plat quasi-compact et quasi-séparé, un $\mathcal{O}_{S}$-Module
quasi-cohérent et plat, et des applications — sont annulées à la machine
lettre par lettre, puis barrées d'un long trait ; elles ne sont plus
lisibles.}

\page{45}
\textbf{11.6.} — Soient $G$ un $S$-groupe plat quasi-compact et
quasi-séparé sur $S$, tel que $\mathcal{A}(G)$ soit un
$\mathcal{O}_{S}$-Module plat $\mathcal{A}$, et $\mathcal{E}$ un
$\mathcal{O}_{S}$-Module quasi-cohérent et plat.
\struck{Nous dirons qu'}\add{Notons qu'}un morphisme de $S$-foncteurs de
$G$ dans
\add{$\operatorname{End}_{\mathcal{O}_{S}}(\mathbb{W}(\mathcal{E}))$}
\struck{\ill{}} $\simeq
\mathcal{H}om_{\mathcal{O}_{S}}(\mathbb{V}(\mathcal{E}), \mathbb{V}(\mathcal{E}))$
définit une \emph{opération} de $G$ sur $\mathcal{E}$, si \add{et seulement
si} le morphisme correspondant
$\theta : G \times_{S} \mathbb{V}(\mathcal{E}) \longrightarrow
\mathbb{V}(\mathcal{E})$ vérifie les deux conditions suivantes :

1\textsuperscript{o}) si $\mu : G \times_{S} G \longrightarrow G$ désigne le
morphisme structural, le diagramme

\begin{tikzcd}[column sep=large]
G \times_{S} G \times_{S} \mathbb{V}(\mathcal{E})
  \arrow[r, "\mathrm{id}_{G} \times \theta"]
  \arrow[d, "\mu \times \mathrm{id}_{\mathbb{V}(\mathcal{E})}"'] &
G \times_{S} \mathbb{V}(\mathcal{E}) \arrow[d, "\theta"] \\
G \times_{S} \mathbb{V}(\mathcal{E}) \arrow[r, "\theta"'] &
\mathbb{V}(\mathcal{E})
\end{tikzcd}

est commutatif.

2\textsuperscript{o}) si $\varepsilon : S \to G$ est la section unité de
$G$, le diagramme

\begin{tikzcd}[column sep=small, row sep=small]
S \times_{S} \mathbb{V}(\mathcal{E})
  \arrow[rr, "\varepsilon \times \mathrm{id}_{\mathbb{V}(\mathcal{E})}"] & &
G \times_{S} \mathbb{V}(\mathcal{E}) \arrow[dl, "\theta"] \\
 & \mathbb{V}(\mathcal{E}) \arrow[ul, "\simeq"] &
\end{tikzcd}

est commutatif.

Il revient au même de dire que le morphisme
$\rho :$ \add{$\mathcal{E} \to$} $\mathcal{A}(G) \otimes_{\mathcal{O}_{S}} \mathcal{E}$
déduit de $\theta$ vérifie les deux axiomes suivants :

\marginal{$\mathcal{E} \to \mathcal{A}(G) \otimes \mathcal{E}$}
\note{La source du morphisme, omise par la machine, est rétablie deux fois
de sa main : au-dessus de la ligne et en marge de gauche.}

(CM1) Si $\Delta : \mathcal{A} \longrightarrow \mathcal{A} \otimes
\mathcal{A} = \mathcal{A}(G \times_{S} G)$ est l'homomorphisme
$\mathcal{A}(\theta)$, le diagramme

\begin{tikzcd}[column sep=large]
\mathcal{A} \otimes \mathcal{A} \otimes \mathcal{A} &
\mathcal{A} \otimes \mathcal{E}
  \arrow[l, "\mathrm{id}_{\mathcal{A}} \otimes \rho"'] \\
\mathcal{A} \otimes \mathcal{E}
  \arrow[u, "\Delta \otimes \mathrm{id}_{\mathcal{E}}"] &
\mathcal{A} \arrow[l, "\rho"'] \arrow[u, "\rho"']
\end{tikzcd}

est commutatif.

\note{Le troisième facteur du sommet supérieur gauche et le sommet
inférieur droit sont écrits $\mathcal{A}$ sur la page, là où l'axiome
demande $\mathcal{E}$ ; la lettre n'est pas distinguée de son $\mathcal{E}$
avec certitude. L'homomorphisme $\Delta$ est celui de $\mu$, et non de
$\theta$ comme le porte la ligne.}

(CM2) Si $\eta : \mathcal{A} \longrightarrow \mathcal{O}_{S}$ est
l'homomorphisme $\mathcal{A}(\varepsilon)$, le diagramme

\begin{tikzcd}[column sep=small, row sep=small]
\mathcal{O}_{S} \otimes \mathcal{E} & &
\mathcal{A} \otimes \mathcal{E}
  \arrow[ll, "\eta \otimes \mathrm{id}_{\mathcal{E}}"'] \\
 & \mathcal{E} \arrow[ul, "\simeq"] \arrow[ur, "\rho"'] &
\end{tikzcd}

est commutatif.

Il revient au même de dire que le morphisme
$\theta_{\mathrm{af}} : G_{\mathrm{af}} \times_{S} \mathbb{V}(\mathcal{E})
\longrightarrow \mathbb{V}(\mathcal{E})$ correspondant à $\rho$ définit une
opération de $G_{\mathrm{af}}$ sur $\mathcal{E}$ \struck{\ill{}}
\add{(comparer I 4.2)}.

Par conséquent, il y a correspondance biunivoque entre opérations de $G$
sur $\mathcal{E}$ et opérations de $G_{\mathrm{af}}$ sur $\mathcal{E}$.

\textbf{Lemme 11.7.} — Soient $G$ un $S$-groupe quasi-compact et
quasi-séparé sur $S$, \add{plat et} tel que $\mathcal{A}(G)$ soit

\page{46}
un $\mathcal{O}_{S}$-Module plat $\mathcal{A}$, et $\mathcal{E}$ un
$\mathcal{O}_{S}$-Module quasi-cohérent. On suppose soit que $G$ est affine
\add{sur $S$}, soit que \struck{$G$} et $\mathcal{E}$ \struck{soit} plats
sur $S$. Soit $\rho : \mathcal{E} \to \mathcal{A} \otimes \mathcal{E}$ un
homomorphisme définissant une opération de $G$ sur $\mathcal{E}$. Soit
$\mathcal{E}_{0}$ un sous-$\mathcal{O}_{S}$-Module quasi-cohérent de
$\mathcal{E}$ tel que l'homomorphisme
$\mathcal{E}_{0} \to \mathcal{A} \otimes \mathcal{E}$ déduit de $\rho$ se
factorise à travers l'homomorphisme canonique \add{(injectif, \ill{} on
fait que $\mathcal{A}$ est plat !)}
$\mathcal{A} \otimes \mathcal{E}_{0} \longrightarrow
\mathcal{A} \otimes \mathcal{E}$ au moyen de l'homomorphisme
$\rho_{0} : \mathcal{E}_{0} \longrightarrow
\mathcal{A} \otimes \mathcal{E}_{0}$. Alors $\rho_{0}$ définit une
opération de $G$ sur $\mathcal{E}_{0}$, qu'on appellera \emph{opération
induite} sur $\mathcal{E}_{0}$ par $\rho$, et on dira que
$\mathcal{E}_{0}$ est \emph{stable} sous $\rho$.
\note{L'incise entre parenthèses est de sa main ; le mot qui l'ouvre après
« injectif, » n'est pas lu.}

Cela résulte immédiatement des définitions et de (11.6). On remarquera
cependant qu'en général l'application canonique
$\mathbb{W}(\mathcal{E}_{0}) \longrightarrow \mathbb{W}(\mathcal{E})$
n'est pas injective.

\textbf{Lemme 11.8.} — Soient $G$ un $S$-groupe quasi-séparé et
quasi-compact sur $S$ tel que $\mathcal{A}(G)$ soit un
$\mathcal{O}_{S}$-Module \struck{\ill{}} libre $\mathcal{A}$ de base
$(\varphi_{i})$, $\mathcal{E}$ un $\mathcal{O}_{S}$-Module quasi-cohérent.
On suppose soit que $G$ est affine \add{sur $S$}, soit que $G$ et
$\mathcal{E}$ sont plats sur $S$. Soit
$\rho : \mathcal{E} \to \mathcal{A} \otimes \mathcal{E}$ un homomorphisme
définissant une opération de $G$ sur $\mathcal{E}$, et soit
$x \in \Gamma(S, \mathcal{E})$. Posons
$\rho(x) = \sum \varphi_{i} \otimes x_{i}$, où les
$x_{i} \in \Gamma(S, \mathcal{E})$ sont bien déterminés par cette
relation). Alors le sous-module \add{$\mathcal{F}$} de $\mathcal{E}$
engendré par les $x_{i}$ est stable sous $\rho$, quasi-cohérent et de type
fini sur $\mathcal{O}_{S}$, et $x \in \Gamma(S, \mathcal{F})$.
\note{La parenthèse fermante après « cette relation » est sans ouvrante ;
telle quelle. La fin de l'énoncé — « quasi-cohérent et de type fini… » —
est retapée au-dessus d'un membre de phrase biffé.}

\struck{Soit $\mathcal{F}$ le sous-module de $\mathcal{E}$ engendré par les
$x_{i}$.} Il est clair que $\mathcal{F}$ est quasi-cohérent et de type
fini. L'axiome (CM1) montre que
$\sum \Delta(\varphi_{i}) \otimes x_{i}
= \sum \varphi_{i} \otimes \rho(x_{i})$ ; d'autre part,
$\Delta(\varphi_{i}) = \sum \varphi_{j} \otimes a_{ij}$, d'où
$\sum \varphi_{j} \otimes a_{ij} \otimes x_{i}
= \sum \varphi_{i} \otimes \rho(x_{i})$, d'où, puisque les $\varphi_{i}$
forment une base de $\mathcal{A}$,
$\rho(x_{i}) = \sum a_{ji} \otimes x_{j}$, où
$a_{ji} \in \Gamma(S, \mathcal{A})$, ce qui montre que
$\rho(\mathcal{F}) \subset \mathcal{A} \otimes \mathcal{F}$, donc
$\mathcal{F}$ est stable sous $\rho$. Enfin l'axiome (CM2) montre que
$x = \sum_{i} \eta(\varphi_{i})\, x_{i} \in \Gamma(S, \mathcal{F})$.
\note{Dans $\rho(x_{i}) = \sum a_{ji} \otimes x_{j}$, l'indice du dernier
$x$ est surfrappé sur la page ; c'est $j$ que la sommation demande.}

\textbf{Proposition 11.9.} — Soit $G$ un $S$-groupe quasi-compact
quasi-séparé sur $S$ tel que $\mathcal{A} = \mathcal{A}(G)$ devienne
localement libre après extension fidèlement plate quasi-compacte de la base
$S$ (ce qui est le cas \add{par exemple} lorsque $G$ est un groupe
réductif, comme nous le verrons dans \struck{\ill{}} XXV), et soit
$\mathcal{E}$ un $\mathcal{O}_{S}$-\struck{Module} quasi-cohérent. On
suppose soit que $G$ est affine \add{sur $S$}, soit que $G$ et
$\mathcal{E}$ sont plats sur $S$. Soit
$\rho : \mathcal{E} \longrightarrow \mathcal{A} \otimes \mathcal{E}$ un
homomorphisme définissant une opération de $G$ sur $\mathcal{E}$.

\page{47}
Alors, pour tout sous-$\mathcal{O}_{S}$-Module quasi-cohérent
$\mathcal{F}$ de $\mathcal{E}$ il existe un plus petit
sous-$\mathcal{O}_{S}$-Module quasi-cohérent $\overline{\mathcal{F}}$ de
$\mathcal{E}$ contenant $\mathcal{F}$ qui soit stable sous $\rho$. Lorsque
$\mathcal{F}$ est de type \struck{\ill{}} fini, il en est de même de
$\overline{\mathcal{F}}$.

\marginal{Sa formation commute à toute extension de la base $S' \to S$}

On se ramène aisément \add{par descente fidèlement plate} au cas où
$\mathcal{A}$ est un $\mathcal{O}_{S}$-Module libre, auquel cas la
proposition est conséquence du lemme 11.8.

\textbf{Corollaire 11.10.} — \struck{Soit}\add{Soient} $S$ un
\struck{pré}schéma \struck{affine}, $G$ un $S$-groupe quasi-séparé
\add{et} quasi-compact sur $S$, tel que $\mathcal{A}(G)$ devienne
\struck{\ill{}} localement libre après extension fidèlement plate
quasi-compacte de la base $S$, et soit $\mathcal{E}$ un
$\mathcal{O}_{S}$-Module quasi-cohérent. On suppose soit que $G$ est
affine sur $S$, soit que $G$ et $\mathcal{E}$ sont plats sur $S$. \add{On
se donne une opération de $G$ sur $\mathcal{E}$.} Alors $\mathcal{E}$ est
limite inductive d'une famille filtrante croissante de
sous-$\mathcal{O}_{S}$-Modules quasi-cohérents de type fini de
$\mathcal{E}$ stables sous $G$.
\note{Une main a récrit au-dessus de la ligne « quasi-compact et
quasi-séparé », dans l'autre ordre.}

C'est une conséquence immédiate de (11.9) car $\mathcal{E}$ est
\struck{engendré par ses}\add{limite inductive de ses} sections (EGA I
1.4.1) \struck{et car l'ensemble des sous-$\mathcal{O}_{S}$-Modules
quasi-cohérents de type fini de $\mathcal{E}$ stables sous $G$ est stable
par engendrement.} \add{sous-Modules quasi-cohérents de type fini (EGA I
9.4.9., IV 1.7.7.)}

\textbf{Remarque 11.10.1.} J.P. Serre a remarqué que dans 11.9. (donc
aussi dans 11.10) il suffisait de supposer que pour toute ouvert affine
$U$ de $S$, il existe un morphisme fidèlement plat $U' \to U$, avec $U'$
affine d'anneau $A'$, tel que $\mathcal{E} \otimes_{U} U'$ soit défini par
un $A'$-module projectif (pas nécessairement libre). De plus, lorsque $S$
est localement noéthérien régulier et de dimension 1, cette hypothèse peut
être remplacée par la seule hypothèse que $G$ soit plat sur $S$ (ce qui
implique que $\mathcal{A}(G)$ est plat) ; voir à ce sujet \struck{un
exposé de Serre dans SGA \add{6} (1966/67).} \add{\textsuperscript{1}
J.-P. Serre, Sur les Groupes de Grothendieck des schémas en groupes
réductifs déployés, Pub. Math. n\textsuperscript{o} \uncertain{34},
p. 37--52.}

\textbf{Proposition 11.11.} — Soient $k$ un corps, $G$ un \struck{corps}
$k$-groupe algébrique. Pour que $G$ soit affine, il faut et il suffit
qu'il existe un espace vectoriel de dimension finie $V$, et une immersion
fermée de $G$ dans $\underline{\mathrm{GL}}_{k}(V)$ qui soit un morphisme
de $k$-groupes.

S'il existe un espace vectoriel $V$ de dimension finie et une immersion
fermée $u : G \to \underline{\mathrm{GL}}_{k}(V)$ puisque
$\underline{\mathrm{GL}}_{k}(V)$ est affine (Exp. I, 4.5), il en est de
même de $G$ (EGA I 4.2.3).

Réciproquement, supposons $G$ affine. Notons $A$ l'algèbre de $G$. Puisque
$G$ opère sur lui-même par translations à gauche, il opère aussi sur $A$,
cette opération se traduisant par l'application \struck{diagonal}
diagonale $\Delta : A \longrightarrow A \otimes_{k} A$. D'après (11.10),
$A$ (en tant qu'espace vectoriel) est réunion d'une famille filtrante
croissante $(V_{i})_{i \in I}$ de sous-espaces vectoriels de dimension
finie stables sous $G$. \struck{\ill{}}

\note{De la Proposition 11.11 au bas de la page, le texte est annulé par
de longues diagonales : l'énoncé est repris page 48 sous une autre forme,
à cinq conditions équivalentes. Les deux dernières lignes sont en outre
annulées à la machine lettre par lettre et ne sont plus lisibles.}

\page{48}
\textbf{Proposition 11.11.} — Soit \struck{un $k$ un corps}, $G$ un
\struck{schéma en} groupe\struck{s} \struck{quasi-compact}
\add{algébrique sur le} \struck{sur} \add{corps} $k$. Les conditions
suivantes sont équivalentes :
\begin{itemize}
\item[(i)] $G$ est affine.
\item[(ii)] $G$ est quasi-affine.
\item[(iii)] On peut faire opérer $G$ fidèlement sur un \add{$k$-}schéma
quasi-affine.
\item[(iv)] On peut faire opérer $G$ fidèlement sur un vectoriel sur $k$
(pas nécessairement de dimension finie).
\item[(v)] \struck{\ill{}} $G$ est isomorphe à un sous-groupe fermé d'un
groupe $\mathrm{Gl}(n)_{k}$.
\end{itemize}
\note{L'énoncé corrigé se lit : « Soit $G$ un groupe algébrique sur le
corps $k$. » Le membre biffé qui ouvre (v) est illisible sous la
rature.}

\emph{Démonstration.} On a (i) $\Rightarrow$ (ii) trivialement, et (ii)
$\Rightarrow$ (iii), car $G$ opère fidèlement sur lui-même par
translations à gauche. On a (iii) $\Rightarrow$ (iv) car si $G$ opère
\struck{\ill{}} sur $X$ quasi-compact et quasi-séparé sur $k$, il opère
aussi sur le vectoriel $\Gamma(X, \underline{\mathcal{O}}_{X})$, grâce au
fait que la formation de $f_{*}(\underline{\mathcal{O}}_{X})$
($f : X \to \operatorname{Spec} k$ le morphisme structural) commute à tout
changement de base $S \to \operatorname{Spec} k$ (EGA IV 1.7.12) ; donc il
opère sur l'enveloppe affine
$X' = \operatorname{Spec}(\Gamma(X, \underline{\mathcal{O}}_{X}))$, et le
morphisme canonique $X \to X'$ est évidemment compatible avec l'action de
$G$ ; ceci dit, si $X$ est quasi-affine i.e. $X \to X'$ est une immersion
ouverte, et si $G$ opère fidèlement sur $X$, il opère fidèlement sur $X'$,
ou ce qui revient au même, sur le vectoriel
$V = \Gamma(X, \underline{\mathcal{O}}_{X})$, ce qui prouve l'implication
(iii) $\Rightarrow$ (iv).
\note{La condition sous laquelle $G$ opère est retapée au-dessus de la
ligne annulée, qui portait « fidèlement sur $X$ quasi-affine sur $k$,
alors $G$ opère fidèlement sur l'enveloppe affine
$X' = \operatorname{Spec}(\dots)$ ».}

Supposons maintenant \struck{\ill{}}\add{(iv), i.e.} que $G$ opère
fidèlement sur le vectoriel $V$. Alors, en vertu de 11.10, $V$ est limite
inductive de sous-espaces vectoriels $V_{i}$ de dimension finie, stables
sous l'action de $G$. Si $K_{i}$ est le noyau de l'action induite de $G$
sur $V_{i}$, i.e. de $G \to \underline{\mathrm{Aut}}_{V_{i}}$, alors
$K_{i}$ est un sous-schéma fermé de $G$, et l'hypothèse que $G$ opère
fidèlement s'exprime aussitôt par le fait que \struck{\ill{}}
l'intersection des $K_{i}$ est le sous-groupe unité de $G$. Comme $G$ est
noethé-

\page{49}
rien, il s'ensuit que l'un des $K_{i}$ est déjà réduit au groupe unité,
donc que $G \to \underline{\mathrm{Aut}}_{V_{i}}$ est un monomorphisme.
C'est donc une immersion fermée en vertu de (1.4.2), ce qui prouve que
(iv) $\Rightarrow$ (v). Comme (v) $\Rightarrow$ (i) trivialement, cela
prouve 11.11 pour $G$ de type fini.

Prenons maintenant le cas général, où il reste à prouver que (iv)
$\Rightarrow$ (i). Avec les notations de la démonstration précédente, soit
$G_{i}$ le \struck{sous-groupe}\add{plus petit} sous-groupe fermé de
$\mathrm{Aut}(V_{i})$ par lequel \add{peut} se factoriser
$G \to \underline{\mathrm{Aut}}_{V_{i}}$ \struck{existe} (un tel groupe
existe, \struck{puisque} comme intersection de tous ceux par lesquels se
factorise $G \to{}$ \struck{\ill{}} $\underline{\mathrm{Aut}}_{V_{i}}$). Si
$V_{j} \supset V_{i}$, on voit alors tout de suite que $G_{j}$ invarie
$V_{i}$, et que le morphisme correspondant
$G_{j} \to \underline{\mathrm{Aut}}_{V_{i}}$ se factorise par $G_{i}$. De
cette façon, les $G_{i}$ forment un système projectif de groupes affines
sur $k$, et on a un morphisme naturel
$G \longrightarrow G' = \varprojlim G_{i}$, où $G'$ est lui-même un groupe
affine (EGA IV 8.2.3). L'hypothèse que $G$ opère \emph{fidèlement} sur $V$
s'exprime alors par le fait que $G \to G'$ est un monomorphisme. On a donc
trouvé un monomorphisme de $G$ dans un schéma affine, ce qui implique que
le schéma quasi-compact \add{et séparé} $G$ est quasi-affine, i.e. que son
faisceau structural est ample, en vertu de EGA IV 18.12.17. En d'autres
termes, posant $H = \operatorname{Spec} \Gamma(G, \underline{\mathcal{O}}_{G})$,
le morphisme canonique $G \to H$ est une immersion ouverte. Notons
maintenant que $H$ est de façon naturelle un schéma en groupes, et
$G \to H$ un homomorphisme de groupes, grâce à 11.2. Cet homomorphisme est
évidemment dominant, il reste à prouver que c'est un isomorphisme, ce qui
résulte \struck{\ill{}} du lemme suivant :

\textbf{Lemme 11.11.1.} — Soient \struck{H}\add{G} un schéma en groupes
sur un corps \add{$k$} \struck{\ill{}}. Alors tout sous-schéma en groupes
ouvert \add{$G$} de $H$ est aussi fermé.

Il revient au même de dire que l'ensemble $H - G$ est ouvert, i.e. est un
voisinage de chacun de ses points $x$. Or, toute extension de corps $k'/k$
définissant un morphisme $\operatorname{Spec} k' \to \operatorname{Spec} k$
qui est universellement

\note{La page est traversée de deux longues diagonales qui se croisent ;
le texte reste lisible d'un bout à l'autre et la démonstration se poursuit
page 50.}

\page{50}
ouvert (EGA IV 2.4.9), il \struck{nous} suffit de prouver qu'il existe une
telle extension $k'$ et un point $x'$ de $H' = H \otimes_{k} k'$ au dessus
de $x$, tels que $H' - G'$ soit un voisinage de $x'$. Or \struck{il}
suffira de prendre \struck{\ill{}} $k' = k(x)$, de sorte que $G'$ admet un
point $x'$ rationnel sur $k'$ au dessus de $x$. Alors l'ouvert dans $H'$
translaté de $G'$ par $x'$ contient $x'$ et est contenu dans $H' - G'$,
\struck{\ill{}} qui est donc un voisinage de $x'$, ce qui achève la
démonstration de 11.11.

\textbf{Remarque 11.11.1.} \struck{Monsieur} M. Raynaud nous signale qu'on
peut généraliser 11.11 au cas d'un schéma en groupes quasi-compact et
quasi-séparé $G$ sur un schéma \add{$S$ localement noethérien} régulier de
dimension $\leqslant 1$, en supposant de plus $G$ \emph{plat} sur $S$. On a
alors équivalence des conditions suivantes :
\begin{itemize}
\item[(i)] $G$ est affine sur $S$.
\item[(ii)] $G$ est quasi-affine sur $S$.
\item[(iii)] On peut faire opérer $G$ fidèlement sur un $S$-schéma
quasi-affine.
\item[(iv)] On peut faire opérer $G$ fidèlement sur un Module
quasi-cohérent plat sur $S$.
\item[(v)] (Si $G$ est de type fini sur $S$ \add{et $S$ noethérien}) $G$
est isomorphe à un sous-groupe fermé d'un $\underline{\mathrm{Aut}}_{V}$,
où $V$ est un Module localement libre de type fini sur $S$.
\end{itemize}

\textbf{Remarque 11.11.2.} Avec les notations de la démonstration
précédente, on voit facilement (une fois qu'on sait que $G$ est affine) que
l'homomorphisme $G \to G' = \varprojlim G_{i}$ est un isomorphisme, ce qui
donne pour $G$ un théorème de structure, qui sera \struck{repris}
\add{prouvé} par une autre méthode dans 11.13.

\note{Deux croix, de la même main que celles de la page 49, annulent ici le
premier paragraphe et la Remarque 11.11.2 ; le texte reste lisible. Un long
crochet de marge relie les deux Remarques.}

\page{51}
Notons $\varphi_{i} : G \longrightarrow \underline{\mathrm{GL}}_{k}(V_{i})$
le morphisme déterminé par l'opération de $G$ sur $V_{i}$, et
$\varphi : G \to \underline{\mathrm{GL}}_{k}(A)$ le morphisme déterminé par
l'opération de $G$ sur $A$. Alors, puisque $A$ est réunion des $V_{i}$,
$\operatorname{Ker} \varphi$ est \uncertain{intrésection} des
$\operatorname{Ker} \varphi_{i}$ ; or $\operatorname{Ker} \varphi$ est le
$k$-groupe unité ; puisque $G$ est noethérien, il existe $i \in I$ tel que
$\operatorname{Ker} \varphi_{i}$ soit le $k$-groupe unité, ce qui implique
(I 1.4.2) que $\varphi_{i}$ est une immersion fermée, c.q.f.d.
\note{« intrésection » est tapé ainsi. Ce paragraphe est annulé par les
mêmes diagonales que les pages 49 et 50 ; il achève la démonstration du cas
général, que la Remarque 11.11.2 renvoie à 11.13.}

\textbf{Lemme 11.12.} — Soient $k$ un corps, $G$ un $k$-groupe
quasi-compact. Posons $A = \mathcal{A}(G)$. Etant donné $x \in A$, il
existe \struck{un sous-$k$-espace vectoriel de dimension finie} une
sous-$k$-algèbre de type fini $V$ de $A$ tel que $x \in V$, que
$\Delta(V) \subset V \otimes_{k} V$ et $u(V) \subset V$, où $u$ désigne
l'involution de $A$ correspondant à la symétrie de $G$.

D'après (11.3), on peut remplacer $G$ par $G_{\mathrm{af}}$, donc on peut
supposer $G$ affine.

Soit $(\varphi_{i})$ une base de $A$. Alors
$\Delta(x) = \sum \varphi_{i} \otimes a_{i}$ et
$\Delta(\varphi_{i}) = \sum \varphi_{j} \otimes b_{ji}$ ; l'axiome (HA1) de
(Exp. I 4.2) montre que
$\sum \varphi_{i} \otimes \Delta(a_{i})
= \sum \Delta(\varphi_{i}) \otimes a_{i}$, d'où
$\Delta(a_{i}) = \sum b_{ij} \otimes a_{j}$.

Soit alors $V$ la \struck{sous-espace vectoriel} sous-$k$-algèbre de $A$
engendrée par les $b_{ij}$ et les $u(b_{ij})$. Il est clair que $V$ est une
$k$-algèbre de type fini. L'axiome (HA1) montre aussi que
\[
\sum \Delta(\varphi_{j}) \otimes b_{ji}
= \sum \varphi_{j} \otimes \Delta(b_{ji}) ,
\]
et on en déduit que $\Delta(b_{ij}) = \sum b_{ik} \otimes b_{kj}$ et que
$\Delta(u(b_{ij})) = \sum u(b_{kj}) \otimes u(b_{ik})$. Puisque $\Delta$
est un homomorphisme d'algèbres, on en déduit que
$\Delta(V) \subset V \otimes_{k} V$ et $u(V) \subset V$. Enfin l'axiome
(HA2) de (Exp. I 4.2) montre que $a_{i} = \sum \eta(a_{j})\, b_{ij}$ et que
$x = \sum \eta(\varphi_{i})\, a_{i}$, si bien que $x \in V$.

\textbf{Proposition 11.13.} — \struck{\ill{}} Soient $k$ un corps et $G$ un
$k$-groupe affine d'algèbre $A$. Alors $G$ est limite projective d'un
système filtrant croissant de $k$-groupes affines de type fini, dont les
morphismes de transition sont fidèlement plats.

Comme l'ensemble des sous-$k$-algèbres de type fini de $A$ stables par
$\Delta$ et $u$ est stable par engendrement, $A$ est limite inductive d'une
famille filtrante croissante $(B_{i})_{i \in I}$ de sous-$k$-algèbres de
type fini stables par $\Delta$ et $u$. Alors, \uncertain{cahque} algèbre
$B_{i}$, munie de l'homomorphisme
$B_{i} \longrightarrow B_{i} \otimes_{k} B_{i}$ déduit de $\Delta$ et de la
restriction de $u$ à $B_{i}$ est munie d'une structure d'hyperalgèbre
associative augmentée involutive, donc (Exp. I 4.2) est la $k$-algèbre d'un
$k$-groupe affine $G_{i}$, de type fini sur $k$. Enfin, puisque
$A = \varinjlim V_{i}$,
\note{Les deux corrections de cette page — « une sous-$k$-algèbre de type
fini » pour « un sous-$k$-espace vectoriel de dimension finie », et
« sous-$k$-algèbre » pour « sous-espace vectoriel » — sont retapées à la
machine au-dessus des membres annulés. Dans la dernière ligne, la limite
est celle des $B_{i}$ ; la page écrit $V_{i}$.}

\page{52}
$G = \varprojlim G_{i}$ (EGA IV 8.2.3). \struck{\ill{}} les morphismes de
transition sont \uncertain{fidèlemnet} plats d'après le lemme suivant :

\textbf{Lemme 11.14.} — Soient $k$ un corps, $G$ et $H$ deux $k$-groupes
affines de type fini, $u : G \to H$ un morphisme de $k$-groupes, et
$u^{\mathrm{o}} : B \to A$ l'homomorphisme de $k$-algèbres correspondant.
Pour que $u$ soit fidèlement plat, il faut et il suffit que
$u^{\mathrm{o}}$ soit injectif.

la condition est évidemment nécessaire (EGA IV 2.2.3 et
$\mathrm{O}_{\mathrm{IV}}$ 6.6.1). Montrons qu'elle est suffisante. Posons
$K = \operatorname{Ker} u$. Alors $G/K$ est un $k$-groupe \add{de type
fini} (Exp. $\mathrm{VI}_{\mathrm{A}}$ 5.2), et $u$ se factorise en
$G \xrightarrow{\ p\ } G/K \xrightarrow{\ v\ } H$, où $p$ est fidèlement
plat (Exp. $\mathrm{VI}_{\mathrm{A}}$ 3.2) et où $v$ est un monomorphisme
de $k$-groupes de type fini, donc une immersion fermée (1.4.2). Puisque $H$
est un schéma affine, et que $v$ est une immersion fermée, $G/K$ est un
schéma affine (EGA I 4.2.3), et, si on note $C$ l'algèbre de $G/K$,
l'homomorphisme $v^{\mathrm{o}} : B \to C$ est surjectif (loc. cit.) Or,
puisque $u^{\mathrm{o}}$ est injectif, et que
$u^{\mathrm{o}} = p^{\mathrm{o}} \circ v^{\mathrm{o}}$, $v^{\mathrm{o}}$
est aussi injectif : c'est un isomorphisme, donc $v$ est un isomorphisme,
et, puisque $p$ est fidèlement plat, il en est de même de $u$.

\textbf{Définitions 11.15.} — Etant donné un corps $k$, un $k$-espace
vectoriel $V$, un $k$-groupe $G$ opérant sur $V$, et un vecteur
$v \in V$, on dit que $v$ est \emph{semi-invariant} sous $G$ si le
sous-espace vectoriel \add{$k.v$} \struck{est} stable sous $G$. Notons $A$
l'algèbre de $G$, et $\rho : V \to A \otimes_{k} V$ l'application linéaire
qui définit l'opération de $G$ sur $V$. Dire que $v$ est semi-invariant
sous $G$ revient à dire \struck{\ill{}} qu'il existe $\lambda \in A$ tel
que, pour tout $w \in kv$, $\rho(w) = \lambda \otimes w$. Lorsque
$v \neq 0$, l'élément $\lambda \in A$ est bien déterminé par cette
relation, on l'appellera \struck{invariant}\add{poids} de l'élément
semi-invariant $v$.

\struck{D'autre part, $v$ s'interprète comme un élément du bidual de $V$,
donc comme une section de $\mathbb{V}(V)$ ; notons
$\varphi : G \times \mathbb{V}(V) \to \mathbb{V}(V)$ le morphisme
définissant l'opération de $G$ sur $V$.} Dire que $v$ est semi-invariant
\struck{\ill{}} d'invariant $\lambda$ revient à dire que
$\varphi(g, v) = \lambda(g)\, v$ pour tout $g \in G(k)$, où
\struck{\ill{}} \add{$\lambda : G \to \underline{\mathbb{G}}_{m,k}$} est le
morphisme défini par $\lambda$. On voit donc que $\lambda$
\struck{correspond à}\add{est} un caractère \struck{$\chi$} de $G$, appelé
\emph{caractère associé} à l'élément semi-invariant $v$. \struck{Enfin nous
appellerons poids de l'élément semi-invariant $v$ le poids du caractère
$\chi$ associé à $v$.}
\note{« fidèlemnet » est tapé ainsi. La dernière phrase, qui définissait le
\emph{poids}, est biffée parce que le mot a été reporté plus haut, sur
« invariant ».}

\page{53}
\textbf{Théorème 11.16.} — (Chevalley) Soient $k$ un corps, $G$ un
$k$-groupe algébrique affine, $H$ un sous-préschéma en groupes fermé de
$G$. Alors il existe un nombre fini d'éléments $a_{i}$ linéairement
indépendants de l'algèbre $A$ de $G$, tels que $H$ soit le plus grand
sous-préschéma en groupes fermé de $G$ sous lequel les $a_{i}$ soient
semi-invariants ; on peut de plus supposer que tous les $a_{i}$ ont même
\struck{caractère associé (donc même} poids\struck{)} sous $H$.

Rappelons que, puisque $G$ est affine, il en est de même de $H$ (Exp. I
4.2.3). Soit $A$ (resp. $B$) l'algèbre de $G$ (resp. $H$) ; d'après (EGA I
4.2.3, il existe un idéal $I$ de $A$ tel que $B$ soit isomorphe à $A/I$.
Puisque $G$ opère sur lui-même par translations à gauche, il opère aussi
sur $A$, cette opération se traduisant par l'application diagonale
$\Delta : A \longrightarrow A \otimes_{k} A$ (Exp. I 4.2 et 4.7.2).
Puisque $G$ est de type fini sur $k$, $A$ est de type fini sur $k$, donc
$I$ admet un système fini de générateurs $(g_{1}, \dots, g_{r})$. Il
résulte de (11.9) qu'il existe un sous-espace vectoriel $V$ de dimension
finie de $A$ \struck{\ill{}} stable sous $G$, et contenant chacun des
$g_{i}$.

Posons alors $W = V \cap I$, c'est un espace vectoriel sur $k$ de dimension
finie, dont nous noterons $d$ la dimension. Puisque $V$ contient tous les
$g_{i}$, $W$ engendre \struck{l'algèbre}\add{l'idéal} $I$.

Notons $p : A \longrightarrow A/I = B$ l'application canonique, et
$q = p \otimes \mathrm{id}_{A} : A \otimes_{k} A \longrightarrow
(A/I) \otimes_{k} A$. Alors l'action de $H$ sur $A$ est déterminée par
$q \circ \Delta : A \longrightarrow B \otimes_{k} A$. Puisque $H$ est un
sous-préschéma en groupes de $G$, l'application diagonale
\[
A/I \to (A/I) \otimes_{k} (A/I)
= ((A/I) \otimes_{k} A) / ((A/I) \otimes_{k} I)
\]
se déduit de $q \circ \Delta$ par passage au quotient, ce qui montre que
$(q \circ \Delta)(I) \subset (A/I) \otimes_{k} I$, autrement dit, $I$ est
stable sous $H$. Puisque $V$ est stable sous $G$, donc sous $H$, on voit
que $W$ est stable sous $H$.

Posons $E = \bigwedge^{d} V$, soit $(w_{1}, \dots, w_{d})$ une base de $W$
sur $k$, et soit $(e_{o}, \dots, e_{n})$ une base de $E$ sur $k$ contenant
$e_{o} = w_{1} \wedge \dots \wedge w_{d}$. L'opération de $G$ sur $V$
détermine canoniquement une opération de $G$ sur $\bigwedge^{d} V$, qui
définit une application linéaire $\rho : E \to A \otimes_{k} E$. Posons
$\rho(e_{i}) = \sum a^{i}_{j} \otimes e_{j}$, et
$a_{i} = a^{o}_{i}$. Puisque $W$ est stable sous $H$, si on pose
$\sigma = (p \otimes \mathrm{id}_{E}) \circ \rho : E \longrightarrow
(A/I) \otimes_{k} E$, il est clair que $\sigma(e_{o})$ est proportionnel à
$e_{o}$, autrement dit, pour $1 \leqslant i \leqslant n$, on a
$a_{i} \in I$. L'axiome (CM1) de (Exp. I 4.7.2) appliqué à

\page{54}
$\rho(e_{o}) = \sum a_{i} \otimes e_{i}$ montre que
$\sum \Delta(a_{i}) \otimes e_{i} = \sum a_{i} \otimes \rho(e_{i})
= \sum a_{i} \otimes a^{i}_{j} \otimes e_{j}$, d'où on déduit que
$\Delta(a_{i}) = a_{o} \otimes a_{i} + \sum_{j \geqslant 1}
a_{j} \otimes a^{j}_{i}$ ; donc, puisque pour $j \geqslant 1$,
$a_{j} \in I$, on voit que
$q \circ \Delta(a_{i}) = p(a_{o}) \otimes a_{i}$, si bien que pour
$1 \leqslant i \leqslant n$, $a_{i}$ est semi-invariant sous $H$
\struck{\ill{}} \struck{d'invariant}\add{de poids} $p(a_{o})$, indépendant
de $i$ ; \struck{donc tous les $a_{i}$ sont semi-invariants sous $H$ de
même caractère associé, donc de même poids. Enfin} \add{Extrayons} du
système $(a_{1}, \dots, a_{n})$, \struck{on peut extraire} un système
maximal de vecteurs linéairement indépendants, dont on peut supposer qu'il
se note $(a_{1}, \dots, a_{m})$. \add{Alors les $a_{i}$, $i = 1, \dots m$
sont linéairement indépendants, semi-invariants sous $H$, et de même
poids.}

Réciproquement, soit $H'$ un sous-préschéma en groupes fermé de $G$ sous
lequel chacun des $a_{i}$ ($1 \leqslant i \leqslant m$) est semi-invariant.
\add{Montrons que $H' = H$.} Il existe de même un idéal $I'$ de $A$ tel que
l'algèbre $B'$ de $H'$ soit isomorphe à $A/I'$ ; notons
$p' : A \to A/I'$ l'application canonique, et
$q' = p' \otimes \mathrm{id}_{A}$.

Par hypothèse, pour $1 \leqslant i \leqslant m$, il existe
$\lambda_{i} \in B'$ tel que
$q' \circ \Delta(a_{i}) = \sum p'(a_{j}) \otimes a^{j}_{i}
= \lambda_{i} \otimes a_{i}$.

Or l'axiome (CM2) de (Exp. I 4.7.2) montre que
$e_{i} = \sum \eta(a^{i}_{j})\, e_{j}$, donc que
$\eta(a^{i}_{j}) = \delta_{ij}$ (symbole de Kronecker) ; l'égalité
$\sum p'(a_{j}) \otimes \eta(a^{i}_{j}) = \lambda_{i} \otimes \eta(a_{i})$
s'écrit donc $p'(a_{i}) = 0$ pour $1 \leqslant i \leqslant m$, car
$\eta(a_{i}) = \eta(a^{o}_{i}) = 0$ pour $1 \leqslant i \leqslant m$ ; par
conséquent, on a $a_{i} \in I'$ pour $1 \leqslant i \leqslant m$, si bien
que $a_{i} \in I'$ pour $1 \leqslant i \leqslant n$. Donc, si on pose
$\sigma' = (p' \otimes \mathrm{id}_{E}) \circ \Delta$,
$\sigma'(e_{o}) = \sum p'(a_{i}) \otimes e_{i}$ est proportionnel à
$e_{o}$, de sorte que
$q' \circ \Delta(W) \subset A/I' \otimes_{k} W$, autrement dit que $W$ est
stable sous $H'$. Puisque $W$ engendre \add{l'idéal} $I$, $I$ est stable
sous $H'$, donc
$q' \circ \Delta(I) \subset (A/I') \otimes_{k} I$, et
$q' \circ \Delta$ définit par passage au quotient une application linéaire
\[
A/I \longrightarrow (A/I') \otimes_{k} (A/I)
= ((A/I') \otimes_{k} A) / ((A/I') \otimes_{k} I)
\]
donc \struck{\ill{}} la restriction à $H' \times H$ du morphisme de
multiplication de $G$ se factorise à travers $H$, ce qui implique que $H'$
est \struck{\ill{}} un sous-préschéma en groupes de $H$, c.q.f.d.
\note{La ligne « la restriction à $H' \times H$… » est retapée
au-dessus d'une ligne annulée à la machine et devenue illisible.}

\textbf{Théorème 11.17 (Chevalley).} — Soient $k$ un corps, $G$ un
$k$-groupe affine (non nécessairement algébrique), et $H$ un
sous-préschéma en groupes fermé de $G$ \emph{invariant} dans $G$ ; alors le
faisceau quotient (pour la topologie fpqc) $G/H$ est représentable par un
$k$-groupe \emph{affine}.

\marginal{Il doit exister une référence ici. Exp. $\mathrm{VI}_{\mathrm{A}}$ \ill{}}
\note{Note au crayon, écrite dans la marge de gauche le feuillet tourné ;
une seconde ligne, plus pâle, n'est pas lue.}

\emph{Premier cas} : Supposons d'abord $G$ algébrique. D'après (Exp.
$\mathrm{VI}_{\mathrm{A}}$ 3.2), le conoyau $K$ de l'injection
$H \to G$ existe, \struck{et} le morphisme canonique
$p : G \longrightarrow K$ est \struck{fidèlement plat, et quasi-compact,
puisque $G$ est de type fini sur $k$, et que $K$, qui est muni d'une
structure de $k$-groupe (Exp. $\mathrm{VI}_{\mathrm{A}}$ 5.2) est séparé
sur $k$ (Exp. $\mathrm{VI}_{\mathrm{A}}$ 0.2). Donc $p$ est} couvrant pour
la topologie fpqc, si bien \struck{(Exp. IV 4.4.3)} \add{et} \struck{que}
$K$ représente le faisceau quotient $G/H$ pour la topologie fpqc ;

\page{55}
\note{Feuillet entièrement de sa main, au crayon : c'est le « \textbf{VOIR CI
CONTRE} » que la page 56 porte entouré, à la place du renvoi biffé à
Chevalley. Il donne le cas où $k$ est algébriquement clos.}

Supposons d'abord $k$ algébriquement clos, $G$ réduit et connexe et $H$
réduit. D'après (\emph{BIBLE} \struck{\ill{}} \add{Exp 4, cor. du
th. 1}), il existe une représentation linéaire de dimension finie
$\rho : G \longrightarrow \mathrm{GL}(V)$ telle que
$(\operatorname{Ker} \rho)_{\mathrm{red}} = H$. Par 11.11,
$G/\operatorname{Ker} \rho$ est affine, donc
$(G/H)/(\operatorname{Ker} \rho/H)$ est affine, donc aussi $G/H$, le groupe
$\operatorname{Ker} \rho/H$ étant fini.

Si $k$ est algébriquement clos, et $G$ et $H$ réduits, alors
$(G/H)/(G^{\mathrm{o}}/H \cap G^{\mathrm{o}})$ est fini, \add{comme
quotient de $G/G^{\mathrm{o}}$, et}
$G^{\mathrm{o}}/H \cap G^{\mathrm{o}}$ est affine, donc $G/H$ est affine.

\struck{N.B. 1) le th. entraîne que si $G$ est affine algébrique et $H$
invariant, il existe une rep. linéaire de dim. finie
$\rho : G \longrightarrow \mathrm{GL}(V)$ telle que
$H = \operatorname{Ker} \rho$.}

\struck{2) La réf. à BIBLE est un corollaire de 11.16, la démonstration est
facile mais un peu longue.}

\struck{3) Dans la réf. à BIBLE, supposer $G$ connexe, car la dém. de BIBLE
est \ill{} dans le cas général.}
\note{Les trois N.B. sont annulés par un faisceau de diagonales, puis le
bloc entier encadré d'un trait. Le mot qui qualifie la démonstration de
BIBLE au 3) n'est pas lu.}

\page{56}
il s'agit donc de montrer que $K$ est affine.

\struck{Le théorème est connu dans le cas où $k$ est algébriquement clos,
et où $G$ et $H$ sont réduits (Chevalley} \add{Supposons d'abord}
\marginal{\textbf{VOIR CI CONTRE}}
\note{« \textbf{VOIR CI CONTRE} » est écrit en grand et entouré : c'est la
page 55, qui remplace le renvoi biffé.}

Dans le cas où $k$ est algébriquement clos, et $G$ et $H$ quelconques,
\struck{\ill{}}, $(G \times_{k} H)_{\mathrm{red}}$ est isomorphe à
$(G_{\mathrm{red}}) \times_{k} (H_{\mathrm{red}})$, donc
$(G_{\mathrm{red}})/(H_{\mathrm{red}})$ est isomorphe à
$(G/H)_{\mathrm{red}}$. Puisque $(G_{\mathrm{red}})/H_{\mathrm{red}}$ est
affine, il en est de même de $G/H$ (EGA I 5.1.10), car, puisque $G/H$ est
de type fini sur $k$, son nilradical est nilpotent.

Dans le cas où $k$ est quelconque,
$G \otimes_{k} \bar{k} \,/\, H \otimes_{k} \bar{k}$ est isomorphe à
$(G/H) \otimes_{k} \bar{k}$ (9.2 v)), donc puisque le premier est affine,
il en est de même du second, donc $G/H$ est affine \add{(réf.)}.

\emph{Deuxième cas} : Cas général. D'après (11.13), l'algèbre $A$ de $G$
est limite inductive d'une famille filtrante croissante
$(A_{i})_{i \in I}$ de sous-algèbres de type fini de $A$ stables par
l'application diagonale et l'involution. D'après
(\struck{Exp.}\add{EGA} I 4.2.3), $H$ est affine, et si on note $B$
l'algèbre de $H$, l'homomorphisme $A \longrightarrow B$ déduit de
l'injection $H \longrightarrow G$ est surjectif. Posons
$I = \operatorname{Ker}(A \to B)$, $I_{i} = I \cap A_{i}$ et
$B_{i} = A_{i}/I_{i}$ ; puisque la structure d'hyperalgèbre associative
augmentée involutive de $A$ passe au quotient à travers $I$, on vérifie
aisément que celle de $A_{i}$ passe au quotient à travers $I_{i}$, si bien
que $B_{i}$ est l'algèbre d'un $k$-groupe affine de type fini $H_{i}$ ; il
est clair que $B = \varinjlim B_{i}$. D'après (EGA I 4.2.3), $H_{i}$ est
isomorphe à un sous-préschéma en \uncertain{goupes} fermé de $G_{i}$.
D'après (6.7) le fait que $H$ soit invariant dans $G$ s'exprime en disant
qu'un certain morphisme $H \times_{k} G \longrightarrow G$ se factorise à
travers l'injection $H \to G$, autrement dit que l'homomorphisme
correspondant $A \longrightarrow B \otimes_{k} A$ se factorise à travers
l'homomorphisme canonique $A \longrightarrow B$ ; on vérifie aisément que
l'homomorphisme $A_{i} \longrightarrow B_{i} \otimes_{k} A_{i}$ défini de
manière analogue se factorise à travers $A_{i} \longrightarrow B_{i}$, donc
que $H_{i}$ est invariant dans $G_{i}$. D'après ce qui a été vu
précédemment, $G_{i}$ étant de type fini sur $k$, $G_{i}/H_{i}$ est
représentable par un $k$-groupe affine $K_{i}$, dont nous noterons $C_{i}$
l'algèbre. Posons alors $C = \varinjlim C_{i}$, et soit
$K = \varprojlim K_{i}$ le $k$-préschéma affine d'algèbre $K$ (EGA IV
8.2.3).
\note{« goupes » est tapé ainsi. La dernière ligne écrit « d'algèbre $K$ » ;
c'est $C$ que la construction donne.}

Montrons que $K$ représente $G/H$ ; pour cela, il suffit de vérifier que
\struck{le diagramme} $H \times_{k} G \rightrightarrows G \times_{K} G$
\struck{est cartésien}, et que le morphisme $G \to K$ est couvrant pour la
topologie fpqc. (\struck{Exp.} IV 4.4.3). Or, \struck{chacun des
diagrammes} \add{pour chaque $i$, on a}
$H_{i} \times_{k} G_{i} \rightrightarrows G_{i} \times_{K_{i}} G_{i}$,
\struck{étant cartésien} \add{donc} il en est de même du \add{morphisme}
diagramme obtenu en passant à la limite projective ; enfin chacun des
morphismes
\marginal{est-ce défini ?}
\note{Note au crayon dans la marge de gauche, le feuillet tourné.}

\page{57}
$G_{i} \longrightarrow K_{i}$ est fidèlement plat (9.2 (xi)), autrement dit
$C_{i}$ est fidèlement plat sur $A_{i}$ ; puisque
$A = \varinjlim A_{i}$ et $C = \varinjlim C_{i}$, on en déduit que $C$ est
fidèlement plat sur $A$, \struck{(cf. Bourbaki, Alg. Comm. \ill{})} si bien
que $G \longrightarrow K$ est un morphisme fidèlement plat. Puisque ce
morphisme est affine, il est quasi-compact, donc couvrant pour la topologie
fpqc, c.q.f.d.

\struck{\textbf{Corollaire 11.18.} — Soient $k$ un corps, $G = D(M)$ un
groupe diagonalisable sur $k$ (I 4.4), \struck{H un sous-groupe de} avec
$M$ un $\mathbb{Z}$-module de type fini, $H$ un sous-groupe algébrique de
$G$. Alors $H$ est également diagonalisable.}
\note{Le Corollaire est annulé par une suite de larges boucles au crayon.
C'est la dernière ligne écrite du tapuscrit : le reste du feuillet est
blanc.}

\section*{Catégorie tannakienne de lien un groupe diagonalisable (pages 59
et 60)}

\note{Deux feuillets d'un tapuscrit étranger au précédent — autre papier,
autre frappe — corrigés de sa main. Ce sont les seules pages du lot qui
touchent aux catégories tannakiennes, c'est-à-dire au sujet du carton. La
lettre $\mathcal{M}$, que la machine n'avait pas, est encrée à la main
partout. La page 58, verso blanc de la page 59, où seule transparaît la
frappe de celle-ci, est absente.}

\page{59}
\emph{Catégorie tannakienne} (sur un corps $k$) \emph{de \struck{\ill{}} lien un groupe diagonalisable} $G = D(\Gamma)$
\add{$\Longleftarrow$} ($\mathcal{M}$ semi-simple, et pour tout objet
simple (ou isotypique) $S$ de $\mathcal{M}$, de rang $d(S)$,
$\operatorname{End}(S)$ est de rang $d(S)^{2}$).

Soit $\mathcal{M}$ cette catégorie. Elle est définie à équivalence près
(respectant structures et contraintes) par un élément \add{$\xi$} de
$H^{2}(k, G)$. On a d'autre part
\[
H^{2}(k, G) \simeq \operatorname{Hom}(\Gamma, \mathrm{Br}(k)) ,
\]
donc $\mathcal{M}$ est définie essentiellement par l'homomorphisme
\struck{\ill{}}
\[
\add{\xi}\, u : \Gamma \longrightarrow \mathrm{Br}(k) .
\]

\struck{Pour tout} La catégorie $\mathcal{M}$ est graduée de type $\Gamma$,
les objets homogènes de type $\Gamma$ sont \add{les objets} isotypiques,
donc pour tout $i \in \Gamma$, il y a exactement un objet simple, à
isomorphisme près, $E_{i}$, \struck{de} homogène de degré $i$. (Il semble
que ce soit une deuxième caractérisation des catégories tannakiennes de
lien $D(\Gamma)$, \add{$\Longleftrightarrow$ le produit de deux
isotypiques \ill{} isotypique}) Le degré de cet élément est égal à l'ordre
\add{$d(i)$} de $u(i)$ dans $\mathrm{Br}(k)$. Donc l'anneau
$K(\mathcal{M})$ est donné par
\[
K(\mathcal{M}) \subset \underline{\mathbb{Z}}[\Gamma] ,
\]
\[
\sum n_{i} e^{i} \in K(\mathcal{M}) \iff n_{i} \equiv 0 \ (d(i))
\quad \text{pour tout } i \in \Gamma
\]
La base canonique de $K(\mathcal{M})$ formée par les objets simples étant
formée des $d(i) e^{i}$. Si $M_{i}$ est l'objet simple de \struck{poids}
degré $i$, alors $\operatorname{End}(M_{i})$ est le corps gauche
\struck{\ill{}} de centre $k$ dont l'invariant dans $\mathrm{Br}(k)$ est
$u(i)$.

\struck{En termes d'un système de générateurs (\ill{}$_{i})_{i \in I}$ de
$\Gamma$, \struck{et} (d'un système d'objets \struck{simples $M_{i}$ de
type} isotypiques degré $i$) et des}
\note{Membre annulé d'un trait ondulé ; deux incises tapées au-dessous,
« i.e. $G = \mathbb{G}_{m}^{I}$ » et « i.e. de $\xi$ », sont accolées de sa
main.}

Si $\Gamma = \mathbb{Z}^{(I)}$, la donnée de $u$ revient à la donnée
d'éléments $\xi_{i} \in \mathrm{Br}(k)$. Choisissant pour tout $i$ un objet
isotypique non nul de degré $i$, soit $E_{i}$, et désignant par
$A_{i} = \operatorname{End}(M_{i})$ l'algèbre d'\uncertain{Azumaya}
\add{de classe $\xi_{i}$,} associée, on reconstitue (à équivalence unique
près respectant tout)

\page{60}
$\mathcal{M}$ en termes de la famille des $A_{i}$ \struck{\ill{}} comme
suite. On considère d'abord la catégories sous-additive
\struck{engendrée} ayant comme objets les éléments $\gamma$ de
$\underline{\mathbb{N}}^{(I)}$, ($\underline{\gamma} =
(\gamma_{i})_{i \in I}$ étant interprété comme
$\bigotimes E_{i}^{\gamma_{i}}$), \struck{ou} avec
$\operatorname{Hom}(\underline{\gamma}, \underline{\gamma}') = 0$ si
$\underline{\gamma} \neq \underline{\gamma}'$ et
$= \bigotimes_{i} A_{i}^{\otimes \gamma_{i}}$ si
$\underline{\gamma} = \underline{\gamma}'$, la composition des morphismes
évidentes. On définit de façon évidente le foncteur $\otimes$ (addition de
$\underline{\mathbb{N}}^{(I)}$ sur les objets, produit tensoriel sur les
homomorphismes) et ses contraintes, de façon que l'on trouve une
sous-$\otimes$-catégorie.

\marginal{attention à contrainte de commutativité}

Son enveloppe Karoubienne est une \add{$\otimes$-}catégorie abélienne, qui
s'identifie à la sous-catégorie pleine de $\mathcal{M}$ formée des éléments
dont les composantes homogènes non nulles sont à degrés
$\in \underline{\mathbb{N}}^{(I)}$, \add{et à un foncteur fibre sur
$\bar{k}$.} D'autre part, pour tout $i$, si $A_{i}$ est de
\add{rang}\struck{degré} $d_{i}^{2}$ (donc $E_{i}$ de rang $d_{i}$), soit
$L_{i} = \det E_{i} = \bigwedge^{d_{i}} E_{i}$, (où $E_{i}$ désigne
l'objet de la catégorie construite correspondant à \struck{\ill{}}
$i$-ème élément de la base canonique de
$\underline{\mathbb{Z}}^{(I)}$, soit $l_{i} \dots$), c'est un objet
préinversible (\add{$x \to$} $\otimes L_{i}$ est pleinement fidèle) et en
inversant les $L_{i}$, on trouve (à \struck{équ} équivalence près
respectant tout) la catégorie $\mathcal{M}$.
\note{Dans $\det E_{i}$, les deux lettres qui suivent le signe $=$ sont
surfrappées et reprises à la main ; l'exposant $d_{i}$ est écrit au-dessus.
Le $\bar{k}$ de l'insertion est repris, entouré, dans la marge de gauche.}

Quand $\Gamma$ n'est pas libre, on trouve une construction analogue en
termes d'un système de générateurs \add{convenable}
$(\gamma_{i})_{i \in I}$, mais en « rendant \struck{de plus} égaux à
$1^{x}$ » de plus certains éléments. Ces constructions sont bien adaptées
pour décrire les $\otimes$-foncteurs de $\mathcal{M}$ dans une
\struck{\ill{}} catégorie tannakienne, et pour décrire le changement de
base. (Pour cette dernière raison notamment, il aurait été maladroit dans
le cas général de prendre les $M_{i}$ simples …). Attention, dans la
description des $\otimes$-foncteurs en question, de tenir compte des
contraintes, notamment de celle de commutativité ; il semble qu'il doive
revenir au même d'exiger que le rang est conservé …

\marginal{A expliciter !}

\end{document}
